% GENERATED FILE. Edit canonical lessons, then run npm run generate:books.
% canonical-source-hash: fnv1a64:977f23e6714789cc
% canonical-lessons: ML-C403-pretam, ML-C403-mantrikan, ML-C403-nidhi, ML-C403-sapam, ML-C403-sapikkuka

\chapter{Ghosts, Magic and Treasure}
\label{ch:ml-ghosts-magic-and-treasure}
\hlchaptermodality{403}

\begin{chapteropening}
\textbf{By the end of this chapter:} \emph{I can say a ghost, a magician, a treasure, a curse and to curse.}

Complete the last lesson of chapter 403: I can say a ghost, a magician, a treasure, a curse and to curse.
\end{chapteropening}

\section[\texorpdfstring{\ml{പ്രേതം} (prētaṁ)}{പ്രേതം (prētaṁ)}]{\ml{പ്രേതം} (prētaṁ) — a ghost}

\label{lesson:ML-C403-pretam}



\begin{quote}
Before the new one: say the Malayalam for a hero, then the Malayalam for a moral.
\end{quote}

\subsection*{You'll want to know: \ml{പ്രേതം}}
\textbf{\ml{പ്രേതം}} — \emph{prētaṁ} — \textquotedblleft{}a ghost\textquotedblright{}.

\begin{quote}
\ml{ആ} \ml{പഴയ} \ml{വീട്ടിൽ} \ml{ഒരു} \ml{പ്രേതം} \ml{ഉണ്ടെന്ന്} \ml{ആളുകൾ} \ml{പറയുന്നു}.
\end{quote}

\begin{quote}
\emph{ā paḻaya vīṭṭil oru prētaṁ uṇṭennŭ āḷukaḷ paṟayunnu.}
\end{quote}

\begin{quote}
People say there is a ghost in that old house.
\end{quote}

A story word; people tell ghost stories, \ml{പ്രേതകഥകൾ}, at night.

\begin{cousinweb}
\textbf{\ml{പ്രേതം}} (\emph{prētaṁ}) is borrowed from Sanskrit \emph{preta}, \textquotedblleft{}departed, dead; the spirit of a dead person\textquotedblright{} (Monier-Williams).
\end{cousinweb}

\begin{grammarlens}[title={What you've built}]
One more word for telling what happened.
\end{grammarlens}

\subsection*{Your turn}
\begin{itemize}
\raggedright
  \item \emph{Say it:} \emph{prētaṁ}
  \item \emph{Say it:} \emph{prētaṁ}, once more
  \item \emph{Say it:} \emph{ā paḻaya vīṭṭil oru prētaṁ uṇṭennŭ āḷukaḷ paṟayunnu.} — People say there is a ghost in that old house.
  \item \emph{Recall:} say the Malayalam for a hero, then the Malayalam for a moral, then say \emph{prētaṁ} again
\end{itemize}

\begin{tcolorbox}[breakable,colback=teal!4,colframe=teal!35!black,title={Before you move on}]
What does \ml{പ്രേതം} mean? (\textquotedblleft{}A ghost\textquotedblright{}.) Say it once more.
\end{tcolorbox}
\section[\texorpdfstring{\ml{മാന്ത്രികൻ} (māntrikan)}{മാന്ത്രികൻ (māntrikan)}]{\ml{മാന്ത്രികൻ} (māntrikan) — a magician}

\label{lesson:ML-C403-mantrikan}



\begin{quote}
Before the new one: say the Malayalam for a moral, then the Malayalam for a ghost.
\end{quote}

\subsection*{You'll want to know: \ml{മാന്ത്രികൻ}}
\textbf{\ml{മാന്ത്രികൻ}} — \emph{māntrikan} — \textquotedblleft{}a magician\textquotedblright{}.

\begin{quote}
\ml{മാന്ത്രികൻ} \ml{തൊപ്പിയിൽ} \ml{നിന്ന്} \ml{ഒരു} \ml{മുയലിനെ} \ml{എടുത്തു}.
\end{quote}

\begin{quote}
\emph{māntrikan toppiyil ninnŭ oru muyaline eṭuttu.}
\end{quote}

\begin{quote}
The magician took a rabbit out of his cap.
\end{quote}

Magic is \ml{മാന്ത്രികം} or \ml{ജാലവിദ്യ}.

\begin{cousinweb}
\textbf{\ml{മാന്ത്രികൻ}} (\emph{māntrikan}) is borrowed from Sanskrit \emph{māntrika}, \textquotedblleft{}a reciter of spells, a sorcerer\textquotedblright{} (CDIAL 10050), from \emph{mantra}, \textquotedblleft{}a spell\textquotedblright{}.
\end{cousinweb}

\begin{grammarlens}[title={What you've built}]
One more word for telling what happened.
\end{grammarlens}

\subsection*{Your turn}
\begin{itemize}
\raggedright
  \item \emph{Say it:} \emph{māntrikan}
  \item \emph{Say it:} \emph{māntrikan}, once more
  \item \emph{Say it:} \emph{māntrikan toppiyil ninnŭ oru muyaline eṭuttu.} — The magician took a rabbit out of his cap.
  \item \emph{Recall:} say the Malayalam for a moral, then the Malayalam for a ghost, then say \emph{māntrikan} again
\end{itemize}

\begin{tcolorbox}[breakable,colback=teal!4,colframe=teal!35!black,title={Before you move on}]
What does \ml{മാന്ത്രികൻ} mean? (\textquotedblleft{}A magician\textquotedblright{}.) Say it once more.
\end{tcolorbox}
\section[\texorpdfstring{\ml{നിധി} (nidhi)}{നിധി (nidhi)}]{\ml{നിധി} (nidhi) — a treasure}

\label{lesson:ML-C403-nidhi}



\begin{quote}
Before the new one: say the Malayalam for a ghost, then the Malayalam for a magician.
\end{quote}

\subsection*{You'll want to know: \ml{നിധി}}
\textbf{\ml{നിധി}} — \emph{nidhi} — \textquotedblleft{}a treasure\textquotedblright{}.

\begin{quote}
\ml{കർഷകൻ} \ml{വയലിൽ} \ml{ഒരു} \ml{നിധി} \ml{കണ്ടെത്തി}.
\end{quote}

\begin{quote}
\emph{karṣakan vayalil oru nidhi kaṇṭetti.}
\end{quote}

\begin{quote}
The farmer found a treasure in the field.
\end{quote}

Also a fund of money set aside.

\begin{cousinweb}
\textbf{\ml{നിധി}} (\emph{nidhi}) is borrowed from Sanskrit \emph{nidhi}, \textquotedblleft{}a setting down; a hoard\textquotedblright{} (CDIAL 7207): what was put down and kept.
\end{cousinweb}

\begin{grammarlens}[title={What you've built}]
One more word for telling what happened.
\end{grammarlens}

\subsection*{Your turn}
\begin{itemize}
\raggedright
  \item \emph{Say it:} \emph{nidhi}
  \item \emph{Say it:} \emph{nidhi}, once more
  \item \emph{Say it:} \emph{karṣakan vayalil oru nidhi kaṇṭetti.} — The farmer found a treasure in the field.
  \item \emph{Recall:} say the Malayalam for a ghost, then the Malayalam for a magician, then say \emph{nidhi} again
\end{itemize}

\begin{tcolorbox}[breakable,colback=teal!4,colframe=teal!35!black,title={Before you move on}]
What does \ml{നിധി} mean? (\textquotedblleft{}A treasure\textquotedblright{}.) Say it once more.
\end{tcolorbox}
\section[\texorpdfstring{\ml{ശാപം} (śāpaṁ)}{ശാപം (śāpaṁ)}]{\ml{ശാപം} (śāpaṁ) — a curse}

\label{lesson:ML-C403-sapam}



\begin{quote}
Before the new one: say the Malayalam for a magician, then the Malayalam for a treasure.
\end{quote}

\subsection*{You'll want to know: \ml{ശാപം}}
\textbf{\ml{ശാപം}} — \emph{śāpaṁ} — \textquotedblleft{}a curse\textquotedblright{}.

\begin{quote}
\ml{മാന്ത്രികന്റെ} \ml{ശാപം} \ml{കാരണം} \ml{രാജകുമാരൻ} \ml{ഒരു} \ml{തവളയായി}.
\end{quote}

\begin{quote}
\emph{māntrikanṟe śāpaṁ kāraṇaṁ rājakumāran oru tavaḷayāyi.}
\end{quote}

\begin{quote}
Because of the magician's curse the prince became a frog.
\end{quote}

The opposite of a blessing, \ml{അനുഗ്രഹം}.

\begin{cousinweb}
\textbf{\ml{ശാപം}} (\emph{śāpaṁ}) is borrowed from Sanskrit \emph{śāpa}, \textquotedblleft{}a curse, abuse\textquotedblright{} (CDIAL 12393).
\end{cousinweb}

\begin{grammarlens}[title={What you've built}]
One more word for telling what happened.
\end{grammarlens}

\subsection*{Your turn}
\begin{itemize}
\raggedright
  \item \emph{Say it:} \emph{śāpaṁ}
  \item \emph{Say it:} \emph{śāpaṁ}, once more
  \item \emph{Say it:} \emph{māntrikanṟe śāpaṁ kāraṇaṁ rājakumāran oru tavaḷayāyi.} — Because of the magician's curse the prince became a frog.
  \item \emph{Recall:} say the Malayalam for a magician, then the Malayalam for a treasure, then say \emph{śāpaṁ} again
\end{itemize}

\begin{tcolorbox}[breakable,colback=teal!4,colframe=teal!35!black,title={Before you move on}]
What does \ml{ശാപം} mean? (\textquotedblleft{}A curse\textquotedblright{}.) Say it once more.
\end{tcolorbox}
\section[\texorpdfstring{\ml{ശപിക്കുക} (śapikkuka)}{ശപിക്കുക (śapikkuka)}]{\ml{ശപിക്കുക} (śapikkuka) — to curse}

\label{lesson:ML-C403-sapikkuka}



\begin{quote}
Before the new one: say the Malayalam for a treasure, then the Malayalam for a curse.
\end{quote}

\subsection*{You'll want to know: \ml{ശപിക്കുക}}
\textbf{\ml{ശപിക്കുക}} — \emph{śapikkuka} — \textquotedblleft{}to curse\textquotedblright{}.

\begin{quote}
\ml{ദേഷ്യം} \ml{വന്ന} \ml{മാന്ത്രികൻ} \ml{രാജാവിനെ} \ml{ശപിച്ചു}.
\end{quote}

\begin{quote}
\emph{dēṣyaṁ vanna māntrikan rājāvine śapiccu.}
\end{quote}

\begin{quote}
The angry magician cursed the king.
\end{quote}

Past: \ml{ശപിച്ചു}, \emph{śapiccu}.

The one cursed takes -\ml{എ}.

\begin{cousinweb}
\textbf{\ml{ശപിക്കുക}} (\emph{śapikkuka}) is built on \ml{ശാപം} (\emph{śāpaṁ}), \textquotedblleft{}a curse\textquotedblright{}, borrowed from Sanskrit \emph{śāpa} (CDIAL 12393), with the Malayalam verb ending -\ml{ിക്കുക}.
\end{cousinweb}

\begin{grammarlens}[title={What you've built}]
One more word for telling what happened.
\end{grammarlens}

\subsection*{Your turn}
\begin{itemize}
\raggedright
  \item \emph{Say it:} \emph{śapikkuka}
  \item \emph{Say it:} \emph{śapikkuka}, once more
  \item \emph{Say it:} \emph{dēṣyaṁ vanna māntrikan rājāvine śapiccu.} — The angry magician cursed the king.
  \item \emph{Recall:} say the Malayalam for a treasure, then the Malayalam for a curse, then say \emph{śapikkuka} again
\end{itemize}

\begin{tcolorbox}[breakable,colback=teal!4,colframe=teal!35!black,title={Before you move on}]
What does \ml{ശപിക്കുക} mean? (\textquotedblleft{}To curse\textquotedblright{}.) Say it once more.
\end{tcolorbox}
